% Compile with:  pdflatex thesis.tex   (twice, for the TOC)
\documentclass[11pt,oneside]{report}

\usepackage[a4paper,margin=1in]{geometry}
\usepackage[T1]{fontenc}
\usepackage{lmodern}
\usepackage{amsmath,amssymb,amsthm}
\usepackage{mathtools}
\usepackage{microtype}
\usepackage{enumitem}
\usepackage{graphicx}
\usepackage{tikz}
\usetikzlibrary{calc,arrows.meta,decorations.markings}
\usepackage{xcolor}
\usepackage{listings}
\usepackage{hyperref}
\hypersetup{
  colorlinks=true,
  linkcolor=black,
  citecolor=black,
  urlcolor=blue!50!black,
  pdftitle={Petal Knots and Polygonal Embeddings},
  pdfauthor={Nick Forbes-Young}
}

% ---------- theorem environments ----------
\newtheorem{theorem}{Theorem}[chapter]
\newtheorem{lemma}[theorem]{Lemma}
\newtheorem{proposition}[theorem]{Proposition}
\newtheorem{corollary}[theorem]{Corollary}
\theoremstyle{definition}
\newtheorem{definition}[theorem]{Definition}
\newtheorem{example}[theorem]{Example}
\newtheorem{remark}[theorem]{Remark}

% ---------- math macros ----------
\newcommand{\R}{\mathbb{R}}
\newcommand{\Z}{\mathbb{Z}}
\newcommand{\bbS}{\mathbb{S}}
\newcommand{\PD}{\mathrm{PD}}
\DeclareMathOperator{\Aut}{Aut}
\DeclareMathOperator{\sign}{sign}
\DeclareMathOperator{\Sym}{Sym}
\newcommand{\pd}{\mathrm{pd}}

% ---------- code listings ----------
\definecolor{codegray}{rgb}{0.96,0.96,0.96}
\definecolor{codeblue}{rgb}{0.20,0.30,0.60}
\definecolor{codegreen}{rgb}{0.13,0.42,0.18}
\lstdefinestyle{cpp}{
  language=C++,
  backgroundcolor=\color{codegray},
  basicstyle=\ttfamily\footnotesize,
  keywordstyle=\color{codeblue}\bfseries,
  commentstyle=\color{codegreen}\itshape,
  showstringspaces=false,
  breaklines=true,
  numbers=left,
  numberstyle=\tiny\color{gray},
  frame=single,
  framerule=0pt,
  xleftmargin=10pt,
  framexleftmargin=10pt,
  belowskip=0pt,
  aboveskip=4pt
}
\lstdefinestyle{shell}{
  basicstyle=\ttfamily\footnotesize,
  backgroundcolor=\color{codegray},
  breaklines=true,
  frame=single,
  framerule=0pt,
  xleftmargin=10pt,
  framexleftmargin=10pt
}

% ---------- title ----------
\title{\Huge\bfseries Petal Knots and Polygonal Embeddings\\[0.6em]
\large A Progress Report Toward a\\ Permutation-to-Polygon Pipeline}
\author{Nick Forbes-Young}
\date{June 2026}

\begin{document}

\frontmatter
\pagenumbering{roman}
\maketitle

\thispagestyle{empty}
\vfill
\noindent
\textit{This document is a milestone summary of work in progress. It collects
the knot-theoretic background, software environment, and locally written code
that have been assembled so far to drive a single workflow: starting from a
petal permutation, generating a polygonal embedding of the corresponding petal
knot, suitable for downstream simplification and analysis.}
\vfill

\tableofcontents

\chapter*{Preface}
\addcontentsline{toc}{chapter}{Preface}

The aim of this work is concrete and modest: given a petal permutation $\pi$,
produce a clean 3D polygonal embedding of the petal knot it encodes, so that
downstream tools can simplify and study it. That sentence hides a stack of
ideas. Knot equivalence is invariant under deformation of three-dimensional
embeddings; the petal projection presumes a non-trivial way to read a knot
type out of a single permutation; and the whole construction has to be made
effective by a program that can be run, checked, and built upon.

The contributors to this stack are several. Schumacher and Cantarella's
\emph{Knoodle} library provides the surrounding machinery: simplification of
planar diagrams via the \texttt{PlanarDiagram} class and the \textsc{Reapr}
algorithm, generation of self-avoiding polygons via \texttt{ClisbyTree}, and,
in version 1.0.2, a trio of command-line tools --- \texttt{knoodlesimplify},
\texttt{knoodleidentify}, and \texttt{knoodledraw} --- that reduce, name, and
render a 3D polygon respectively. The locally written code in this directory
adds the petal construction on top of that foundation, together with
sampling and cross-validation tools that use the identification step to
tabulate the knot types produced by whole families of petal permutations.

The chapters that follow lay out the material in the order it is needed.
Chapter~\ref{chap:background} reviews enough knot theory to read the rest of
the document: knots, equivalence, diagrams, Reidemeister moves, and PD codes.
Chapter~\ref{chap:petal} introduces petal knots and petal permutations.
Chapter~\ref{chap:knoodle} describes the parts of Knoodle we depend on.
Chapter~\ref{chap:petalstar} describes the locally written code that
connects a permutation to a polygon. Chapter~\ref{chap:enum} then closes the
loop with the sampling and identification tools: it introduces the compact
knot notation we use, describes three enumeration programs corresponding to
three progressively stronger equivalences on petal permutations, reports the
cross-validation against an independent reference, and records a small
conjecture about constant-difference permutations that emerged along the way.

\mainmatter
\pagenumbering{arabic}

\chapter{Introduction}\label{chap:intro}

A knot, classically, is the image of an embedding $K \colon \bbS^1 \hookrightarrow
\R^3$, considered up to ambient isotopy. The interest of the subject is that
this image is a one-dimensional object floating in three space, with rich
combinatorial shadows. The standard combinatorial shadow is the
\emph{knot diagram}: a planar projection of the knot in which transverse
double points are decorated with over/under information. From a diagram we
can attempt to read off invariants, and from invariants we can attempt to
distinguish knots.

The computational construction we focus on here is straightforward to state.
We have:
\begin{enumerate}[label=(\roman*)]
\item a notion of \emph{petal knot}, encoded by a permutation $\pi$ of $\{1, \dots, n\}$ with $n$ odd;
\item a routine that turns $\pi$ into a polygonal embedding $K(\pi) \subset \R^3$;
\item the Knoodle library, which simplifies the resulting 3D polygon to a
planar diagram code (PD code) and looks it up against a table of known knots
to produce a named knot type.
\end{enumerate}
The full pipeline is therefore
\[
\pi \;\xrightarrow{\;\texttt{petal\_star\_clean}\;}\; K(\pi)
     \;\xrightarrow{\;\texttt{knoodlesimplify}\;}\; \PD(K(\pi))
     \;\xrightarrow{\;\texttt{knoodleidentify}\;}\; \text{knot name}.
\]
The current state of this pipeline is summarized in Chapter~\ref{chap:state}.

\chapter{Knot-Theoretic Background}\label{chap:background}

This chapter develops the language we use in the rest of the document. The
exposition is informal but precise enough to make the algorithms exact. The
standard references are Rolfsen~\cite{Rolfsen} and Kawauchi~\cite{Kawauchi};
for the computational viewpoint, Murasugi~\cite{Murasugi} and Lickorish~\cite{Lickorish}.

\section{Knots, links, and equivalence}

\begin{definition}
A \emph{knot} is a smooth (or piecewise-linear) embedding $K \colon \bbS^1 \to \R^3$.
A \emph{link} of $\ell$ components is a smooth embedding of $\ell$ disjoint
copies of $\bbS^1$. Two knots $K_0, K_1$ are \emph{equivalent} (or
\emph{ambient isotopic}) if there exists a continuous family
$h_t \colon \R^3 \to \R^3$, $t \in [0,1]$, of homeomorphisms with $h_0 = \mathrm{id}$
and $h_1 \circ K_0 = K_1$.
\end{definition}

By a knot $K$ we typically mean its equivalence class. The simplest knot is the
\emph{unknot}, equivalent to the standard circle in a plane. Two knots that
look very different may be equivalent, and one of the central questions of
knot theory is to decide, given two diagrams, whether they represent the same
knot.

\section{Knot diagrams}

A \emph{regular projection} of a knot $K \subset \R^3$ is a projection
$p \colon \R^3 \to \R^2$ whose image $p(K)$ is a closed curve with finitely many
self-intersections, each of which is a transverse double point. A
\emph{knot diagram} is a regular projection augmented with over/under data at
every crossing: at each double point, one of the two strands is marked as
passing \emph{over} the other.

\begin{figure}[h]
\centering
\begin{tikzpicture}[scale=0.8, line width=0.8pt]
% trefoil
\draw[white,line width=4pt] (0,1.4) .. controls (-1.4,1.4) and (-1.4,-1) .. (0,-1)
   .. controls (1.4,-1) and (1.4,0.6) .. (0,0.6);
\draw (0,1.4) .. controls (-1.4,1.4) and (-1.4,-1) .. (0,-1)
   .. controls (1.4,-1) and (1.4,0.6) .. (0,0.6);
\draw[white,line width=4pt] (0,0.6) .. controls (-1.4,0.6) and (-1.4,-0.7) .. (0,-0.7)
   .. controls (1.4,-0.7) and (1.4,1.1) .. (0,1.1);
\draw (0,0.6) .. controls (-1.4,0.6) and (-1.4,-0.7) .. (0,-0.7)
   .. controls (1.4,-0.7) and (1.4,1.1) .. (0,1.1);
\draw[white,line width=4pt] (0,1.1) .. controls (-0.6,1.1) and (-0.6,1.4) .. (0,1.4);
\draw (0,1.1) .. controls (-0.6,1.1) and (-0.6,1.4) .. (0,1.4);
\node at (3.5,0) {trefoil ($3_1$)};
\end{tikzpicture}
\caption{The trefoil knot $3_1$, the simplest nontrivial knot.}
\end{figure}

Two diagrams represent equivalent knots if and only if one can be transformed
into the other by ambient isotopy of the plane and a finite sequence of
local moves called \emph{Reidemeister moves}.

\begin{theorem}[Reidemeister]
Two knot diagrams represent equivalent knots if and only if they are related
by a finite sequence of moves of types I, II, and III (Figure~\ref{fig:reide}),
together with planar isotopies.
\end{theorem}

\begin{figure}[h]
\centering
\begin{tikzpicture}[scale=0.7,line width=0.6pt]
\draw (0,0) -- (1.5,0); \node at (0.75,-0.6) {(I)};
\draw[->] (1.7,0) -- (2.7,0);
\draw (2.9,0) .. controls (3.4,0) and (3.6,0.4) .. (3.8,0)
        .. controls (4,-0.4) and (4.2,0) .. (4.7,0);
\begin{scope}[xshift=6cm]
\draw (0,-0.4) -- (1.5,-0.4); \draw (0,0.4) -- (1.5,0.4);
\node at (0.75,-1.0) {(II)};
\draw[->] (1.7,0) -- (2.7,0);
\draw (2.9,0.4) .. controls (3.6,0.4) and (3.6,-0.4) .. (4.3,-0.4);
\draw[white,line width=3pt] (2.9,-0.4) .. controls (3.6,-0.4) and (3.6,0.4) .. (4.3,0.4);
\draw (2.9,-0.4) .. controls (3.6,-0.4) and (3.6,0.4) .. (4.3,0.4);
\end{scope}
\begin{scope}[xshift=12cm]
\node at (1,-1.0) {(III)};
\draw (0,0) -- (2,0); \draw (0,0.8) -- (2,0.8);
\draw[white,line width=3pt] (0.3,-0.5) -- (1.7,1.3);
\draw (0.3,-0.5) -- (1.7,1.3);
\end{scope}
\end{tikzpicture}
\caption{Schematic Reidemeister moves I, II, III. Each is the local picture of
an allowed transformation between equivalent knot diagrams.}
\label{fig:reide}
\end{figure}

This theorem reduces the topological question of knot equivalence to a purely
combinatorial one. Unfortunately the combinatorics is still hard: the number
of Reidemeister moves needed to relate two equivalent diagrams is not bounded
by any reasonable function of their crossing counts.

\section{PD codes}

To put a knot diagram into a computer we use a \emph{planar diagram code} (PD
code). The standard convention, which Knoodle and our tools follow, is the
one used by the KnotTheory package and by Regina:

\begin{definition}\label{def:pdcode}
Fix an orientation of the knot. Number the arcs of the diagram --- the maximal
segments of the curve between successive under-crossings --- by $1, 2, \dots, n$
in the order they are encountered when traversing the knot once. (Equivalently,
the curve is broken into $2n$ \emph{edges} between consecutive crossing
incidences, labelled $1$ through $2n$.) A PD code for an $n$-crossing diagram
is the list
\[
\PD = \bigl[ X[a_1,b_1,c_1,d_1], X[a_2,b_2,c_2,d_2], \dots, X[a_n,b_n,c_n,d_n] \bigr],
\]
where at each crossing the four edges incident to the crossing are listed in
counter-clockwise order beginning with the incoming under-edge. A
\emph{signed} PD code appends $\pm 1$ to each crossing indicating handedness.
\end{definition}

\begin{example}\label{ex:trefoilpd}
The right-handed trefoil $3_1$ has PD code
\[
\PD = \bigl[ X[1,4,2,5], X[3,6,4,1], X[5,2,6,3] \bigr],
\]
with all three crossings positive. The over-strand at the crossing $X[1,4,2,5]$
runs from edge $4$ to edge $5$ (so edges $4$ and $5$ belong to the same arc);
the under-strand runs from edge $1$ to edge $2$.
\end{example}

\section{Knot invariants}

A \emph{knot invariant} is a function on knot diagrams whose value depends only
on the underlying knot type. The classical invariants include the
\emph{crossing number} (the minimum number of crossings over all diagrams of
$K$), the \emph{Alexander polynomial} $\Delta_K(t) \in \Z[t,t^{-1}]$, the
\emph{Jones polynomial}, and various determinants, signatures, and
homological constructions (Khovanov, Heegaard--Floer). For the present
report we will not need any of these in detail; we only need to know that
they exist and that they are computed from PD codes by Knoodle's downstream
tools.

\chapter{Petal Knots}\label{chap:petal}

The construction of interest in this work is the \emph{petal projection}, due
to Adams \emph{et al.}~\cite{Adams}. A petal projection encodes a knot by a
single permutation. This makes it especially convenient as a parameterization
of an enumeration problem: every knot type appears for some permutation, and
many permutations give the same knot.

\section{Definition}

\begin{definition}
A \emph{petal projection} of a knot is a diagram in the plane consisting of $n$
straight chords (\emph{petals}) emanating from a common central point, joined
into a single closed curve in a prescribed cyclic order. Each pair of petals
crosses exactly once, and all crossings occur at the central point. The
projection is decorated by a height function $h \colon \{1, \dots, n\} \to \R$
that assigns each petal a unique vertical level; the over/under data at the
central crossings is read from this function.
\end{definition}

To resolve the geometric degeneracy of $n$ chords meeting at one point we
shift each chord very slightly off-centre so that they cross pairwise. This is
exactly what the construction in Chapter~\ref{chap:petalstar} does.

\begin{theorem}[Adams, Crawford, DeMeo, et al.~\cite{Adams}]
Every knot admits a petal projection. The minimum number of petals required
is an invariant of the knot, called its \emph{petal number}.
\end{theorem}

\section{Petal permutations}

If we fix the cyclic ordering of the petals (the order in which the closed
curve visits them) and the height function takes distinct integer values
$1, 2, \dots, n$, the petal projection is completely determined by the
permutation $\pi \in \Sym(n)$ where $\pi(t)$ is the height of the $t$-th
petal in traversal order.

\begin{definition}
A \emph{petal permutation} of odd length $n$ is a permutation
$\pi = (\pi_1, \dots, \pi_n)$ of $\{1, \dots, n\}$, where $n$ is odd.
\end{definition}

The oddness condition reflects the standard convention that the traversal step
$s = (n+1)/2$ used to visit petals in turn must be coprime to $n$; this is
automatic for odd $n$ and ensures every petal is visited exactly once.

Two permutations may give the same knot. Operations preserving the knot type
include cyclic shift (which simply re-labels which strand we call ``first''),
reversal (reversal of the orientation), and inversion (relabelling heights).
The composition of these generates a symmetry group of order up to $4n$ acting
on petal permutations.

\begin{example}\label{ex:permops}
For $\pi = (2,4,1,5,3)$ of length $5$:
\begin{itemize}[leftmargin=2em]
\item \textbf{reverse}: $(3,5,1,4,2)$
\item \textbf{inverse}: $\pi^{-1}(j) = i$ where $\pi_i = j$, here $(3,1,5,2,4)$
\item \textbf{cyclic shift by 1}: $(4,1,5,3,2)$
\end{itemize}
\end{example}

In the codebase these operations live in \texttt{PetalPermutation}
(Section~\ref{sec:petalperm}).

\chapter{The Knoodle Library}\label{chap:knoodle}

\emph{Knoodle} is a C++17/20 library for computational knot theory by
Henrik Schumacher and Jason Cantarella. We rely on the \texttt{PlanarDiagram}
class, the \textsc{Reapr} simplification algorithm, and --- as of version
1.0.2 --- the three command-line tools \texttt{knoodlesimplify},
\texttt{knoodleidentify}, and \texttt{knoodledraw}.

\section{Overview}

The library is documented in its own README; we recapitulate only the parts we
use here. The high-level data flow Knoodle is designed to support is
\[
\text{3D polygon (TSV)} \;\to\; \text{planar diagram} \;\to\; \text{(simplified) PD code} \;\to\; \text{invariants}.
\]

\subsection{\texttt{PlanarDiagram} and \texttt{PlanarDiagramComplex}}

The \texttt{PlanarDiagram} class is the central data structure. It stores the
combinatorics of a knot diagram (crossings, arcs, faces, handedness) in a
form amenable to local moves and rewrites. Construction is supported from
several sources, including:
\begin{itemize}[leftmargin=2em]
\item \texttt{FromKnotEmbedding(coords, vertex\_count)}, which takes a 3D polygon
and projects it to obtain a diagram;
\item \texttt{FromSignedPDCode} and \texttt{FromUnsignedPDCode}, which read PD codes
directly.
\end{itemize}
\texttt{PlanarDiagramComplex} extends this with explicit handling of multiple
prime summands, used during simplification when a knot decomposes as a
connect sum.

\subsection{\textsc{Reapr}}

\textsc{Reapr} is Knoodle's flagship diagram-simplification algorithm. It
applies Reidemeister moves and strand reroutings to a diagram, guided by
an energy functional, repeatedly until no further reduction is possible.
It is far more aggressive than the textbook Reidemeister algorithm and
routinely reduces hundred-crossing projection diagrams to their minimal
forms.

For our purposes, \textsc{Reapr} is the engine inside \texttt{knoodletool}
that takes the messy polygon projection of a petal-star embedding and
reduces it to a short, canonical PD code.

\subsection{\texttt{ClisbyTree} and \texttt{PolyFold}}

\texttt{ClisbyTree} is a fast generator of self-avoiding polygons, and
\texttt{PolyFold} is a command-line tool that uses it for random sampling.
We do not currently use either for the petal-knot pipeline, but they are
worth knowing about for any future generalization to random knot
experiments.

\section{The v1.0.2 command-line tools}

Knoodle v1.0.2 replaced the earlier monolithic \texttt{knoodletool} with
three focused command-line tools. All three read 3D polygon coordinates in
the same tab-separated format as produced by \texttt{petal\_star\_clean}.

\subsection{\texttt{knoodlesimplify}}

Runs the \textsc{Reapr} pipeline on the input polygon and prints a
simplified PD code in the convention of Definition~\ref{def:pdcode}. The
output format uses a \texttt{k} line to separate knots and an \texttt{s}
line to separate prime summands; crossings are $5$-tuples of integers
(four edge labels and a sign). A \texttt{--streaming-mode} flag lets many
polygons be simplified in one invocation, which matters when driving it
from a sampling loop.

\begin{lstlisting}[style=shell,caption={Simplify a polygon to a PD code},captionpos=b]
./knoodlesimplify petal_star.tsv > pd_code.txt
\end{lstlisting}

\subsection{\texttt{knoodleidentify}}

Simplifies (via the same \textsc{Reapr} pipeline) and then \emph{names} the
result by looking it up in Knoodle's KLUT (Knot Look-Up Table). The output
is a Wolfram-language association whose keys record the crossing number,
Rolfsen table index, alternating flag, and a slash-separated symmetry
coset:
\begin{lstlisting}[style=shell]
<|KnotSymbol[3,1,True,"e/r"] -> 1|>
\end{lstlisting}
Composite knots come back as a length-$k$ association with one entry per
prime summand. If the input diagram simplifies to the empty diagram, the
tool returns the empty association \texttt{<||>}, which we interpret as
the unknot.

The tool requires the environment variable \texttt{KNOODLE\_KLUT\_DIR}
to point at the KLUT data directory (typically
\texttt{/home/linuxbrew/.linuxbrew/Cellar/knoodle/1.0.2/data/Klut} on
Linuxbrew). Setting it once in \texttt{\textasciitilde/.bashrc} is enough.

\subsection{\texttt{knoodledraw}}

Renders a simplified diagram as ASCII art. It is useful for eyeballing
whether a claimed knot name is plausible, especially when investigating
unfamiliar cases returned by \texttt{knoodleidentify}. Because
\texttt{knoodledraw} draws whatever diagram it is handed, it should
generally be run downstream of \texttt{knoodlesimplify} rather than on the
raw polygon:
\begin{lstlisting}[style=shell]
./knoodlesimplify --streaming-mode < petal_star.tsv | ./knoodledraw
\end{lstlisting}

\chapter{From Permutation to Polygon: the Petal-Star Construction}\label{chap:petalstar}

This chapter walks through the locally written code that converts a petal
permutation into a 3D polygon. There are two implementations,
\texttt{petal\_star\_clean.cpp} and \texttt{petal\_star\_polygon.cpp}.
They differ in elaborateness but pursue the same goal.

\section{The \texttt{PetalPermutation} class}\label{sec:petalperm}

The file \texttt{petal\_permutation.hpp} defines a single class:

\begin{lstlisting}[style=cpp,caption={Selected interface of \texttt{PetalPermutation} (\texttt{petal\_permutation.hpp})},captionpos=b]
class PetalPermutation {
public:
    explicit PetalPermutation(std::vector<int> values);    // validates n odd
    const std::vector<int>& values() const;
    int size() const;
    PetalPermutation reversed() const;
    PetalPermutation inverse() const;
    PetalPermutation rotatedLeft(int shift) const;
    std::vector<PetalPermutation> cyclicShifts() const;
    std::vector<PetalPermutation> symmetryOrbit() const;
    PetalPermutation canonicalRepresentative() const;
    std::vector<int> descents() const, ascents() const,
                     peaks() const, valleys() const;
    std::string toString() const, compactKey() const;
};
\end{lstlisting}

The constructor validates that the input is a permutation of $\{1,\dots,n\}$
with $n$ odd, throwing \texttt{std::invalid\_argument} otherwise. The
\emph{symmetry orbit} of $\pi$ is the set of all permutations obtained by
composing cyclic shifts, reversal, and inversion (the four generators noted
in Example~\ref{ex:permops}). The \emph{canonical representative} is the
lexicographically smallest member of this orbit, providing a unique label
for the orbit.

The combinatorial statistics (\texttt{descents}, \texttt{ascents}, \texttt{peaks},
\texttt{valleys}) record positions in the cyclic sequence where the permutation
falls, rises, or makes a local extremum. These do not affect the polygon
construction, but they are useful for diagnostic output via the
\texttt{--info} flag and for studies of which combinatorial features
correlate with which knot types.

\section{The clean construction: \texttt{petal\_star\_clean.cpp}}

The cleanest implementation is in \texttt{petal\_star\_clean.cpp}. Given a
permutation $\pi$ of length $n$, it builds a 3D polygon consisting of $n$
straight chords through (or just past) the origin, joined by short triangular
``petals'' on the outer rim. The key design choices are:

\begin{enumerate}[leftmargin=2em]
\item \textbf{Traversal step.} The $t$-th chord enters the unit circle at
angle $a_t = 2\pi t s / n$, where $s = (n+1)/2$ is coprime to $n$. Successive
chords therefore exit at adjacent boundary positions, so the connecting
arcs (the ``petals'') are thin and contribute no extra crossings.
\item \textbf{Heights.} Each chord is at the constant height
$z_t = h \cdot \bigl(\pi_t - \tfrac{n+1}{2}\bigr)$, so the vertical order of
the $n$ strands at the centre is exactly the permutation $\pi$.
\item \textbf{Resolving the central degeneracy.} The $n$ chords through the
origin would all cross at a single point (a degenerate $n$-fold crossing).
Each chord's far endpoint is rotated by a uniform small angle
$\delta = \pi/(2n)$; the resulting chords miss the centre and pairwise
intersect cleanly, forming a star of $\binom{n}{2}$ simple crossings.
\end{enumerate}

The over/under data at every central crossing is determined purely by the
strand heights, which in turn is determined purely by $\pi$.

\begin{lstlisting}[style=cpp,caption={Polygon construction (\texttt{petal\_star\_clean.cpp})},captionpos=b]
std::vector<Point3D> buildPetalStar(
    const PetalPermutation& perm,
    double radius, double petalRadius, double heightScale
) {
    const int n = perm.size();
    const auto& vals = perm.values();
    const int step = (n + 1) / 2;
    const double twoPi = 2.0 * M_PI;
    const double delta = M_PI / (2.0 * n);  // uniform shift left

    std::vector<Point3D> polygon;
    for (int t = 0; t < n; ++t) {
        const int    p     = (t * step) % n;
        const int    pNext = ((t + 1) * step) % n;
        const double a     = twoPi * p / n;
        const double aNext = twoPi * pNext / n;
        const double z     = heightFromValue(vals[t], n, heightScale);
        const double zNext = heightFromValue(vals[(t + 1) % n], n, heightScale);

        const double entryAng = a;
        const double exitAng  = a + M_PI + delta;
        const Point3D entry = { radius*cos(entryAng), radius*sin(entryAng), z };
        const Point3D exit  = { radius*cos(exitAng),  radius*sin(exitAng),  z };

        const double gap    = wrapTo2Pi(aNext - exitAng);
        const double tipAng = exitAng + 0.5 * gap;
        const Point3D tip   = { petalRadius*cos(tipAng),
                                petalRadius*sin(tipAng),
                                0.5 * (z + zNext) };
        polygon.push_back(entry);
        polygon.push_back(exit);
        polygon.push_back(tip);
    }
    polygon.push_back(polygon.front());  // close the loop
    return polygon;
}
\end{lstlisting}

The polygon has $3n + 1$ vertices: $n$ chord entries, $n$ chord exits, $n$
petal tips, plus a duplicate of the first vertex to close the loop.

\paragraph{Why $\delta = \pi/(2n)$?} The constraint is $0 < \delta < \pi/n$.
The lower bound is needed to break the degeneracy; the upper bound ensures
that successive chord exits do not overlap the next chord's entry, i.e.\
that the petal gap remains positive. Choosing $\delta = \pi/(2n)$ centres
the chord exit between two extreme positions, which is the simplest choice
that respects both bounds. The resulting diagram is symmetric under
$\pi \mapsto \pi^{-1}$ and $\pi \mapsto \pi^{\mathrm{rev}}$, so no preferred
direction or chirality is introduced by the construction itself.

\paragraph{Verification.} For $n=5$ the polygon has $\binom{5}{2} = 10$ central
crossings, matching the expected count for a $5$-petal projection. The
order of heights at the centre exactly reproduces $\pi$, as required.

\section{The elaborated construction: \texttt{petal\_star\_polygon.cpp}}

\texttt{petal\_star\_polygon.cpp} is an earlier, more elaborate implementation
that constructs the same overall figure but with five named points per strand
($Q_{\mathrm{in}}, P_{\mathrm{in}}, P_{\mathrm{out}}, Q_{\mathrm{out}},
A$) and several layers of perturbation (\texttt{innerJitter},
\texttt{outerJitter}, \texttt{petalJitter}, \texttt{apexLift}). The extra
machinery exists to dodge floating-point degeneracies in earlier projection
code and to keep the projected diagram readable when drawn. The output
polygon serves the same role downstream.

For new development we recommend \texttt{petal\_star\_clean.cpp} as the
canonical generator. Its geometry is simpler to reason about, and its
output has the smallest possible number of vertices while still being a
valid embedding.

\section{Output format}

Both programs write a TSV file (default \texttt{petal\_star.tsv}) containing
one $x,y,z$ triple per line, ten decimal digits of precision. This is the
format consumed by \texttt{knoodletool}.

\chapter{Enumerating and Identifying Petal Knots}\label{chap:enum}

The construction of Chapter~\ref{chap:petalstar} gives one direction of the
pipeline: a permutation goes in, a polygon comes out. Chapter~\ref{chap:knoodle}
gives the other direction: a polygon (through \texttt{knoodleidentify}) yields
a named knot type. Composed, the two directions turn a petal permutation into
a knot name entirely by shell command. This chapter documents the small
collection of sampling programs built on top of that composition, together
with the compact naming convention we use for the results and a first
conjecture that fell out of running them.

\section{Compact knot notation}\label{sec:compact}

\texttt{knoodleidentify} returns a Wolfram-language association whose keys
are of the form
\[
\texttt{KnotSymbol}[c, i, \texttt{alt}, \texttt{"sym"}],
\]
where $c$ is the crossing number, $i$ the Rolfsen table index,
\texttt{alt} \texttt{= True} for alternating knots, and \texttt{sym} a
slash-separated listing of a coset of the diagram's symmetry group (e.g.\
\texttt{e/r}, \texttt{m/mr}, \texttt{e/r/m/mr}). This is precise but
verbose.

For human-readable output we compact each summand to a single string of the
form
\[
\texttt{<chirality><c>.<i>[n]},
\]
with the following convention:
\begin{itemize}[leftmargin=2em]
\item \texttt{chirality} $\in \{\texttt{a}, \texttt{p}, \texttt{m}, \texttt{h}\}$,
where \texttt{a} is amphichiral (coset contains both \texttt{e} and
\texttt{m}), \texttt{p} is ``positive'' (contains only \texttt{e}),
\texttt{m} is mirror (contains only \texttt{m}), and \texttt{h} covers
the rare chiral non-invertible cases whose coset contains neither
(e.g.\ $9_{42}$);
\item \texttt{c.i} is the Rolfsen crossing number and index;
\item the suffix \texttt{n} is appended when $c \geq 11$ and the knot is
non-alternating, matching the convention of the standard tables at
these crossing numbers.
\end{itemize}
Composite knots are joined with \texttt{\#}: for example
\texttt{p3.1\#m3.1} is the connect sum of a positive trefoil and a
negative trefoil. The unknot is \texttt{a0.1}.

\paragraph{Chirality convention.}
Independent implementations of the petal-knot pipeline disagree on which
mirror is ``positive.'' In our reference dataset (see
Section~\ref{sec:crossval}) the \texttt{e}-coset corresponds to
\texttt{m} and the \texttt{m}-coset to \texttt{p} --- exactly the
reverse of what \texttt{knoodleidentify}'s coset naming suggests. We
therefore swap the two branches when computing the chirality prefix.
The swap is documented in-line in the source of every tool that emits
compact knot names, with the original mapping preserved as a comment.

\section{Sampling programs}\label{sec:sampling}

\subsection{\texttt{sample\_all\_perms.cpp}}

\texttt{sample\_all\_perms} enumerates one representative per \emph{rotation
orbit} of petal permutations of length $n$: every permutation with
$\pi_1 = 1$, in lexicographic order. For each rep it invokes
\texttt{petal\_star\_clean} to build the polygon, \texttt{knoodleidentify}
to name the knot, converts the KnotSymbol to compact form, and emits

\begin{lstlisting}[style=shell]
# 1 2 3 4 5 6 7
a0.1

# 1 2 3 4 5 7 6
a0.1
\end{lstlisting}

Cyclic shifts of a petal permutation produce the same knot, so it is
enough to walk the $(n-1)!$ reps with $\pi_1 = 1$ rather than all $n!$
permutations. For $n = 7$ that is $720$ rows; for $n = 9$, $40320$.

\subsection{\texttt{sample\_const\_diff.cpp}}\label{sec:sampleconstdiff}

\texttt{sample\_const\_diff} sweeps the \emph{constant-difference}
permutations
\[
\pi^{(d)}_i = ((i-1)\, d \bmod n) + 1, \qquad \gcd(d, n) = 1,
\]
identifies each resulting knot, and annotates it against a small hardcoded
table of known torus knots (T(2,3), T(2,5), T(2,7), T(2,9), T(3,4),
T(3,5), \ldots). The tool is motivated by the conjecture reported in
Section~\ref{sec:torus-conj}: every constant-difference petal permutation
appears empirically to produce a torus knot. Anything not in the table is
annotated \texttt{?}, so unknowns show up as gaps in a column rather than
as false positives.

\subsection{\texttt{sample\_rotation\_classes.cpp}}\label{sec:samplerotclasses}

\texttt{sample\_rotation\_classes} enumerates canonical representatives of
a \emph{stronger} equivalence than pure rotation: the combined action of
cyclic rotation of position and cyclic shift of value. Two permutations
are identified iff their cyclic-difference tuples
\[
\pd_i(\pi) \;=\; (\pi_{i+1} - \pi_i) \bmod n, \qquad
\pi_{n+1} := \pi_1,
\]
are cyclic rotations of one another. Each orbit has a unique canonical
rep: the permutation with $\pi_1 = 1$ whose $\pd$ tuple is the
lexicographically smallest cyclic rotation of the class. For $n = 9$
this produces $4492$ rows, reproducing the format of an external
reference dataset (\texttt{rclass9.txt}) used for combinatorial
cross-checks.

Rows are further stratified by the invariants
\[
w(\pi) \;=\; \tfrac{1}{n}\sum_i \pd_i(\pi) \qquad \text{(winding number)}
\]
and the \emph{robust} flag, which is true iff every $\pd_i \in \{2, \dots,
n-2\}$ (equivalently, no cyclically adjacent pair of petals has
consecutive heights). Robust classes are emitted first, then non-robust,
each block sorted by $(w, \pd)$.

\section{Equivalences on petal permutations}\label{sec:equivalences}

The three sampling tools of Section~\ref{sec:sampling} correspond to three
progressively stronger equivalences on the space of petal permutations. The
orbit counts under each equivalence give a sense of the scale of the
enumeration.

\begin{center}
\begin{tabular}{@{}lrrr@{}}
\hline
Equivalence & $n=5$ & $n=7$ & $n=9$ \\
\hline
None (all $n!$)                           & 120  & 5040 & 362880 \\
Rotation only ($\pi_1 = 1$)               & 24   & 720  & 40320  \\
Rotation $\times$ value-shift ($C_n \times C_n$) & 8 & 130 & 4492 \\
\hline
\end{tabular}
\end{center}

The rotation-only equivalence is a genuine knot invariant: cyclic shift of
a petal permutation does not change the resulting knot. The rotation
$\times$ value-shift equivalence is coarser and is not, in general, a
knot invariant --- it groups together permutations whose crossing patterns
have the same $\pd$-tuple shape, which is convenient for compact
enumeration and matches the format of external reference data.

Symmetry generators that we do \emph{not} include in either equivalence
(and which would give further identifications) are reversal of the
sequence ($\pi \mapsto \pi^{\mathrm{rev}}$) and value reflection
($\pi_i \mapsto n + 1 - \pi_i$); both produce the mirror image of the
underlying knot.

\section{Cross-validation}\label{sec:crossval}

For $n = 7$, the output of \texttt{sample\_all\_perms} was cross-checked
against an independently produced reference file
(\texttt{Parsley\_n7.txt}) covering the same $720$ rotation classes in
the same lexicographic order. Once the chirality convention of
Section~\ref{sec:compact} is applied, the two agree byte-for-byte after
whitespace normalization: zero knot mismatches across all $720$ rows.
This is strong evidence that the composition
\[
\text{permutation} \longrightarrow
\texttt{petal\_star\_clean} \longrightarrow
\texttt{knoodleidentify}
\]
correctly identifies the knot type of every petal permutation of length
$7$, and by extension provides confidence in the same pipeline at
$n = 9$ and beyond.

The comparison itself is performed by a small utility,
\texttt{file\_compare.cpp}, which strips whitespace from both files and
reports the first non-whitespace byte at which they disagree. This is
robust to differences in trailing newlines, blank-line conventions, and
minor spacing.

\section{The torus-knot conjecture}\label{sec:torus-conj}

An observation that fell out of running \texttt{sample\_const\_diff} is
that every constant-difference petal permutation appears to produce a
torus knot. For $n = 9$ the six values of $d$ coprime to $9$ produce:

\begin{center}
\begin{tabular}{@{}rlll@{}}
\hline
$d$ & permutation & compact knot & torus label \\
\hline
$1$ & $(1,2,3,4,5,6,7,8,9)$ & \texttt{a0.1}       & unknot        \\
$2$ & $(1,3,5,7,9,2,4,6,8)$ & \texttt{m5.1}       & $T(2,5)$      \\
$4$ & $(1,5,9,4,8,3,7,2,6)$ & \texttt{m15.41185n} & (conjectured $T(4,9)$) \\
$5$ & $(1,6,2,7,3,8,4,9,5)$ & \texttt{p15.41185n} & (mirror of above)      \\
$7$ & $(1,8,6,4,2,9,7,5,3)$ & \texttt{p5.1}       & $T(2,5)$      \\
$8$ & $(1,9,8,7,6,5,4,3,2)$ & \texttt{a0.1}       & unknot        \\
\hline
\end{tabular}
\end{center}

The pairing $d \leftrightarrow n - d$ produces mirror-image knots. This is
because replacing $d$ by $n - d$ reverses the cyclic order in which the
petals are visited and so flips the sign of every crossing.

Two directions worth pursuing:
\begin{itemize}[leftmargin=2em]
\item \textbf{Prove the conjecture.} Identify which torus knot
$\pi^{(d)}$ produces as a function of $(n, d)$; the classical
expectation is $T(d, n)$ or a small variant thereof, with the mirror
symmetry above.
\item \textbf{Extend \texttt{TORUS\_TABLE}.} As larger $n$ are swept,
add newly encountered compact names to the recognition table in
\texttt{sample\_const\_diff.cpp} so that the annotation column fills
in without manual look-up.
\end{itemize}

\chapter{Current State and Next Steps}\label{chap:state}

\section{What works}

The current pipeline turns a petal permutation into a named knot type. The
components are:
\begin{itemize}[leftmargin=2em]
\item \texttt{petal\_permutation.hpp}: petal permutation class.
\item \texttt{petal\_star\_clean.cpp}: deterministic permutation-to-polygon
generator. This file is treated as canonical and is not to be modified.
\item \texttt{petal\_star\_polygon.cpp}: earlier, more elaborate polygon
generator, kept for reference.
\item \texttt{sample\_all\_perms.cpp}: enumerates one rep per rotation
orbit of petal permutations of length $n$ and identifies each knot.
\item \texttt{sample\_const\_diff.cpp}: sweeps constant-difference
permutations, identifies each knot, and annotates against a small
table of torus knots.
\item \texttt{sample\_rotation\_classes.cpp}: enumerates canonical reps
under the combined position-rotation $\times$ value-shift equivalence;
reproduces the $4492$-row \texttt{rclass9} reference format for $n = 9$.
\item \texttt{file\_compare.cpp}: whitespace-insensitive file diff, used
for cross-validation against external reference data.
\end{itemize}
The TSV output of \texttt{petal\_star\_clean} feeds
\texttt{knoodlesimplify} and \texttt{knoodleidentify}, giving a
permutation-to-knot-name pipeline that is validated at $n = 7$ against an
independent reference (Section~\ref{sec:crossval}).

An earlier attempt at a home-grown Alexander-polynomial-based identifier
(\texttt{identify\_knot.cpp}) is kept in the repository for reference but
has been superseded by \texttt{knoodleidentify}, which is both faster and
handles a much wider range of knots via the KLUT.

\section{Open directions}

\begin{itemize}[leftmargin=2em]
\item \textbf{The torus-knot conjecture.} Prove that every
constant-difference petal permutation produces a torus knot and pin
down the exact $T(p, q)$ as a function of $(n, d)$
(Section~\ref{sec:torus-conj}). Extending the torus recognition
table in \texttt{sample\_const\_diff.cpp} at larger $n$ is the
immediate empirical step.
\item \textbf{Direct \texttt{perm\_to\_pd}.} For petal knots specifically,
the PD code is determined analytically by the permutation: there are
exactly $\binom{n}{2}$ central crossings, the over/under data is read
from the permutation, and the traversal order is fixed by the step
$s = (n+1)/2$. A direct \texttt{perm\_to\_pd} routine would remove the
polygon-projection round-trip from the front of the pipeline.
\item \textbf{Scaling to $n = 11$ and beyond.} At $n = 11$ the rotation
orbit count reaches $10! = 3.6$M, which is on the edge of feasibility
for the current shell-out driver. A streaming or in-process
integration with \texttt{knoodleidentify} (avoiding process startup
per permutation) is the natural performance improvement.
\item \textbf{Cataloguing under further symmetries.} The
\texttt{PetalPermutation} class already computes symmetry orbits under
reversal and inversion; these can drive an even coarser enumeration
by identifying mirror-related knots when that is the appropriate
level of granularity.
\end{itemize}

\begin{thebibliography}{9}

\bibitem{Adams}
C.~Adams, T.~Crawford, B.~DeMeo, et al.,
\emph{Knot projections with a single multi-crossing},
J.\ Knot Theory Ramifications \textbf{24} (2015), no.~3.

\bibitem{Kawauchi}
A.~Kawauchi,
\emph{A Survey of Knot Theory},
Birkh\"auser, 1996.

\bibitem{Knoodle}
H.~Schumacher and J.~Cantarella,
\emph{Knoodle: a library for computational knot theory},
\url{https://github.com/HenrikSchumacher/Knoodle}.

\bibitem{Lickorish}
W.~B.~R.~Lickorish,
\emph{An Introduction to Knot Theory},
Springer GTM 175, 1997.

\bibitem{Murasugi}
K.~Murasugi,
\emph{Knot Theory and Its Applications},
Birkh\"auser, 1996.

\bibitem{Rolfsen}
D.~Rolfsen,
\emph{Knots and Links},
AMS Chelsea, 2003.

\end{thebibliography}

\end{document}
